\section{Referências, citações e siglas}

\subsection{Referências cruzadas}

Associe um nome ao elemento com \verb|\label| e recupere seu número com
\verb|\ref|. A página é obtida com \verb|\pageref|. Por exemplo, a
Seção~\ref{sec:estrutura} apresentou a estrutura do projeto. A
Figura~\ref{fig:unesp} está na página~\pageref{fig:unesp}, a
Figura~\ref{fig:subfiguras} contém duas imagens, a Tabela~\ref{tab:exemplo}
contém dados tabulares e a Equação~\ref{eq:energia} mostra uma expressão
numerada. O símbolo \verb|~| impede uma quebra de linha entre o nome e o número.

\begin{lstlisting}[language={[LaTeX]TeX},caption={Referência cruzada}]
Veja a Figura~\ref{fig:unesp} na pagina~\pageref{fig:unesp}.
\end{lstlisting}

\subsection{Citações bibliográficas}

Cadastre cada obra em \texttt{referencias.bib}. Use \verb|\textcite| quando o
nome fizer parte da frase, \verb|\parencite| para uma citação entre parênteses e
\verb|\cites| para mais de uma obra. Exemplos: \textcite{dummy2024article},
\parencite{dummy2024book} e \cites{dummy2024conference}{dummy2024thesis}.

\begin{lstlisting}[language={[LaTeX]TeX},caption={Comandos de citação}]
Segundo \textcite{dummy2024article}, o tema e relevante.
Outros trabalhos tratam do assunto \parencite{dummy2024book}.
Compare tambem \cites{dummy2024conference}{dummy2024thesis}.
\end{lstlisting}

\subsection{Siglas}

O comando \verb|\nomenclature| registra uma sigla na lista pré-textual. Ele pode
ser colocado na primeira ocorrência do termo.

\begin{lstlisting}[language={[LaTeX]TeX},caption={Cadastro de uma sigla}]
UNESP\nomenclature{UNESP}{Universidade Estadual Paulista}
\end{lstlisting}

UNESP\nomenclature{UNESP}{Universidade Estadual Paulista}

\subsection{Código-fonte}

Use \texttt{lstlisting} para trechos maiores e \verb|\verb| para um comando curto
no meio do parágrafo. A opção \texttt{language} seleciona o realce de sintaxe;
\texttt{caption} cria a legenda e \texttt{label} permite referências, como no
Código~\ref{cod:estrutura}.

\begin{lstlisting}[language=Python,caption={Exemplo de código Python},label={cod:python}]
def saudacao(nome):
    return f"Ola, {nome}!"

print(saudacao("UNESP"))
\end{lstlisting}

O Código~\ref{cod:python} será incluído automaticamente na lista de códigos.

\subsection{Links}

Use \verb|\url{https://www.latex-project.org}| para exibir um endereço e
\verb|\href{https://www.latex-project.org}{Projeto LaTeX}| para escolher o texto
do link. Resultado: \href{https://www.latex-project.org}{Projeto \LaTeX}.
